\documentclass[11pt,reqno]{amsart}

%==============================================================
% Beyond Correlation: Operator-Derived Features for
%   Compositional Risk
%   A Machine-Learning Hedge Experiment with an Iterated-Brownian
%   Prototype
% --------------------------------------------------------------
% Companion to:
%   Fontes, Operator Clark--Ocone Formula from Primitive
%     Divergence  [OperatorClarkOcone]
%   Fontes, Energy-First Stochastic Calculus  [EnergyFirst]
%   Fontes, First-Chaos Dual Closure of the IBM Ito Correction
%     [FCDC]
%   Fontes, A Bidegree Atlas for the IBM Ito Correction
%     [BidegreeAtlas]
%   Fontes, From Carre-du-Champ Hierarchies to Primitive
%     Divergence  [GammaHierarchyDivergence]
%
% Scope: finance-facing conceptual framework paper. Compositional
% randomness X_t = Phi(\Xi_{tau_t}) as a distinct architecture from
% correlational multi-factor models. IBM Z_t = B(Y_t) as the
% canonical mathematically tractable prototype. Operator
% derivatives D^B (field/scenery) and D^Y (clock/state) separate
% the two risk directions; mollified product-source Clark--Ocone
% decomposes hedge demand. The unmollified clock-derivative
% singularity is the mathematical signature of rough compositional
% risk. No no-arbitrage, no calibration, no martingale
% representation in the natural IBM observable filtration.
%==============================================================

\usepackage[utf8]{inputenc}
\usepackage[T1]{fontenc}
\usepackage{lmodern}
\usepackage[margin=1.15in]{geometry}
\usepackage{amsmath,amssymb,amsthm,mathtools}
\usepackage{mathrsfs}
\usepackage{microtype}
\usepackage{hyperref}
\usepackage{xurl}
\usepackage{cleveref}
\usepackage{graphicx}
\usepackage{booktabs}
\usepackage{array}
\usepackage{tabularx}
\usepackage{makecell}
\usepackage{threeparttable}

\emergencystretch=3em

% --- theorem environments
\theoremstyle{plain}
\newtheorem{theorem}{Theorem}[section]
\newtheorem{proposition}[theorem]{Proposition}
\newtheorem{lemma}[theorem]{Lemma}
\newtheorem{corollary}[theorem]{Corollary}

\theoremstyle{definition}
\newtheorem{definition}[theorem]{Definition}
\newtheorem{example}[theorem]{Example}
\newtheorem{assumption}[theorem]{Assumption}
\newtheorem{convention}[theorem]{Convention}
\newtheorem{hypothesis}[theorem]{Hypothesis}

\theoremstyle{remark}
\newtheorem{remark}[theorem]{Remark}

% --- macros
\newcommand{\R}{\mathbb R}
\newcommand{\N}{\mathbb N}
\newcommand{\Z}{\mathbb Z}
\newcommand{\E}{\mathbb E}
\newcommand{\Prob}{\mathbb P}
\newcommand{\Dom}{\mathrm{Dom}}
\newcommand{\Ran}{\mathrm{Ran}}
\newcommand{\cA}{\mathcal A}
\newcommand{\cD}{\mathcal D}
\newcommand{\cE}{\mathcal E}
\newcommand{\cF}{\mathcal F}
\newcommand{\cN}{\mathcal N}
\newcommand{\cU}{\mathcal U}
\newcommand{\IBM}{Z}
\newcommand{\Bvar}{B}
\newcommand{\Yclock}{Y}
\newcommand{\UGam}{\cU_{\Gamma}}
\newcommand{\Piad}{\Pi^{\mathrm{ad}}}
\newcommand{\Lzero}{L^{2}_{0}}
\DeclareMathOperator{\supp}{supp}
\DeclareMathOperator{\Var}{Var}
\DeclareMathOperator{\Cov}{Cov}

\title[Operator-Derived Features for Compositional Risk]{Beyond
       Correlation: Operator-Derived Features for
       Compositional Risk\\
       \normalsize A Machine-Learning Hedge Experiment with an
       Iterated-Brownian Prototype}
\author{Ramiro Fontes}
\address{Quijotic Research}
\email{ramirofontes@gmail.com}
\urladdr{www.quijotic.com}
\thanks{Code, hyperparameters, and CSV outputs supporting the
numerical experiment of \cref{sec:experiment} are available at
\url{https://github.com/QuijoticResearch/compositional-risk-features}.}
\date{\today}

\begin{document}

\begin{abstract}
On synthetic compositional data, augmenting a fixed
non-compositional baseline with six operator-derivative
features raises out-of-sample $R^{2}$ by approximately
$0.5$--$0.7$ across four payoffs (softplus, quadratic, tanh,
smoothed call). Ridge regression on the operator features
essentially reproduces the analytic first-order operator hedge;
gradient-boosted trees add value in the large-clock-shock
regime. On a correlation-matched control with the same marginal
variance, the gap collapses to $|\Delta R^{2}|\le 0.01$,
isolating the advantage to compositional structure.

Many asset observables are compositions
$X_{t}=\Phi(\Xi_{\tau_{t}})$: a random market field $\Xi$
sampled along a random clock/state process $\tau$. The
filtration, topology, and field/clock-sensitivity content of
this composition are not visible to covariance/PCA descriptions
unless the compositional architecture is modeled explicitly.
The operator-derivative framework
of~\cite{GammaHierarchyDivergence,OperatorClarkOcone} supplies
two directional operator derivatives---a field/scenery
derivative $D^{\Bvar}$ and a clock/state derivative
$D^{\Yclock}$---and from these a small set of state-conditioned
compositional features. Iterated Brownian
motion $\IBM_{t}=\Bvar(\Yclock_{t})$ is the mathematical
prototype: the cleanest setting in which the architecture is
exposed.

The framework's role is feature engineering for compositional
risk: operator derivatives identify the features, ML does the
function approximation.
\end{abstract}

\maketitle

%==============================================================
\section{Introduction}\label{sec:intro}
%==============================================================

The standard quantitative-finance literature models market risk
as a \emph{correlational} structure: a vector of latent factors
$W^{(1)},\ldots,W^{(k)}$ drives observable processes via linear
combinations or smooth transforms,
\begin{equation}\label{eq:corr-model}
  dX^{(i)}_{t}\;=\;\mu_{i}(t,X_{t})\,dt
   \;+\;\sum_{j=1}^{k}\sigma_{ij}(t,X_{t})\,dW^{(j)}_{t},
   \qquad d\langle W^{(i)},W^{(j)}\rangle_{t}=\rho_{ij}(t,X_{t})\,dt.
\end{equation}
The covariance matrix $(\rho_{ij})$ glues the drivers and
parametrizes co-movement; portfolio risk decomposes into factor
sensitivities, with the matrix as the central structural input
for co-movement (alongside drift, the diffusion surface, the
state variables, and the filtration).

Many real market observables are not naturally of this form.
They are \emph{compositions}:
\begin{equation}\label{eq:comp-model}
  X_{t}\;=\;\Phi\bigl(\Xi_{\tau_{t}}\bigr),
\end{equation}
where $\Xi$ is a random market field and $\tau_{t}$ is a random
clock / state / sampling process at which the field is read.
Examples developed in \cref{sec:finance} include latent
value/liquidity/risk-premium surfaces sampled by order-flow,
volatility surfaces sampled by random liquidity regime,
execution-cost fields sampled by random trade intensity, and
credit/liquidity-state observables sampled by event clocks. The
composition $\Phi(\Xi_{\tau_{t}})$ is not a linear combination
of independent drivers, and is not adequately described by a
covariance matrix. The conceptual punchline is short:
\emph{field risk} is sensitivity to the random market surface
$\Xi$ at a given sampling location; \emph{clock/state risk} is
sensitivity to moving through that surface, i.e.\ to changes in
the sampling process $\tau$; correlation-only models collapse
this distinction into a single off-diagonal covariance entry.

The minimal mathematically tractable example is \emph{iterated
Brownian motion} (IBM),
\begin{equation}\label{eq:IBM}
  \IBM_{t}\;=\;\Bvar(\Yclock_{t}),\qquad t\ge 0,
\end{equation}
with $\Bvar=(\Bvar_{x})_{x\in\R}$ a two-sided Brownian motion
(scenery / random field) and $\Yclock=(\Yclock_{s})_{s\ge 0}$ an
independent Brownian motion (clock / state). IBM is a
well-studied compositional Gaussian process; for path properties
and quartic-variation behavior see Burdzy~\cite{Burdzy}. It is
used here as a mathematical prototype, not a calibrated price
model: its non-semimartingality, fourth-order PDE shadow,
$L^{2}$-divergent It\^o correction, and quartic operator-energy
dual~\cite{EnergyFirst,FCDC,BidegreeAtlas} all stem from the
compositional structure rather than from any high-dimensional
or non-Gaussian feature.

The practical test of the framework is operational: at any
maturity $T$ and any payoff $F=\Phi(X_{T})$, does ML given
compositional features outperform the same ML given only
non-compositional (covariance/PCA-based) features, on
compositional data? The numerical experiment of
\cref{sec:experiment} answers yes by
$\Delta R^{2}\approx 0.5$--$0.7$ across four payoffs, with a
covariance-matched negative control collapsing the gap. The
contribution is feature construction, not replacement of ML:
the operator derivative identifies a small set of
state-conditioned compositional features that a standard ML
hedger (gradient-boosted regression trees) cannot recover from
covariance/PCA features alone.

\subsection{Operator derivatives identify the features}\label{sec:intro-operators}

The operator-derivative framework
of~\cite{GammaHierarchyDivergence,OperatorClarkOcone} gives two
directional operator derivatives on a payoff
$F=f(\Bvar(\Yclock_{T}))$, on the product source space
$(\Omega_{\Bvar}\times\Omega_{\Yclock},\Prob_{\Bvar}\otimes\Prob_{\Yclock})$:
the field/scenery derivative $D^{\Bvar}_{x}F$ measures the
response of the payoff to perturbations of the realized field
value $\Bvar(x)$ at coordinate $x$; the clock/state derivative
$D^{\Yclock}_{s}F$ measures the response to perturbations of the
clock path $\Yclock(s)$. These are not in general the classical
Malliavin derivatives with respect to the underlying white
noises (\cref{sec:DB-DY}). The operator Clark--Ocone
formula~\cite[\S 7]{OperatorClarkOcone} decomposes the centered
payoff into field and clock components. Reading the chain rule for the
cylindrical IBM payoff $F=f(\Bvar_{\varepsilon}(\Yclock_{T}))$
(\cref{sec:operator-derivatives,sec:carre}) identifies six
state-conditioned compositional features:
\[
  \delta\Bvar_{\varepsilon}(\Yclock),\quad
  \partial_{x}\Bvar_{\varepsilon}(\Yclock)\,\delta\Yclock,\quad
  f'(X)\,\delta\Bvar_{\varepsilon}(\Yclock),\quad
  f'(X)\,\partial_{x}\Bvar_{\varepsilon}(\Yclock)\,\delta\Yclock,\quad
  \Gamma_{\mathrm{field}},\quad\Gamma_{\mathrm{clock}},
\]
where $X=\Bvar_{\varepsilon}(\Yclock)$ and
$\Gamma_{\mathrm{field}}$, $\Gamma_{\mathrm{clock}}$ are the
squared carr\'e-du-champ terms of \cref{sec:carre}. These six
features are the object of the experiment in
\cref{sec:experiment}.

\subsection{Rough compositional risk}\label{sec:intro-rough}

For the cylindrical IBM payoff $F=f(\Bvar(\Yclock_{T}))$ the
clock derivative $D^{\Yclock}_{s}F$ involves the spatial
derivative $\partial_{x}\Bvar(\Yclock_{T})$, which is not an
ordinary $L^{2}$-random variable since $\partial_{x}\Bvar$ is
white noise in space. After spatial mollification
$\Bvar\mapsto\Bvar_{\varepsilon}=\rho_{\varepsilon}*\Bvar$ the
clock derivative is well-defined; the unmollified limit lives
in a strictly weaker operator-energy module identified
in~\cite{EnergyFirst,FCDC,BidegreeAtlas}. The mollified
derivatives are used throughout; the unmollified limit is not
treated here. This singularity is the mathematical signature of
\emph{rough compositional risk}: when a random field is sampled
at a stochastic state, the clock/state-hedge demand can fail to
be classical $L^{2}$, and identifying the correct
operator-energy topology becomes the central analytic problem.

\medskip

\Cref{sec:correl-vs-comp} contrasts correlational and
compositional risk architectures.
\Cref{sec:IBM-prototype} sets up the IBM prototype.
\Cref{sec:operator-derivatives} develops the operator-derivative
decomposition into field and clock components.
\Cref{sec:carre} introduces the compositional carr\'e-du-champ
and the role of mollification, motivating the squared
$\Gamma$-features. \Cref{sec:operator-hedging} states the
analytic operator hedge via the operator Clark--Ocone formula,
kept brief, as it motivates the feature set rather than serving
as a standalone hedge method. \Cref{sec:finance} discusses
finance-facing interpretations of the field/clock split.
\Cref{sec:experiment} reports the numerical experiment. The
operator-derivative machinery used here---closable
$\Gamma^{(2)}$-seeds, primitive divergence, an operator
Clark--Ocone formula, and the IBM trilogy of energy-first
calculus, first-chaos dual closure, and the bidegree
atlas~\cite{GammaHierarchyDivergence,OperatorClarkOcone,EnergyFirst,FCDC,BidegreeAtlas}---is
part of a broader operator-derivative reformulation; the
present paper is one practical instantiation of that machinery
for the finance audience.

%==============================================================
\section{Correlational versus compositional randomness}\label{sec:correl-vs-comp}
%==============================================================

Fix a probability space $(\Omega,\cF,\Prob)$ with a filtration
$(\cF_{t})_{t\ge 0}$ and a finite vector of independent
$(\cF_{t})$-Brownian motions $W^{(1)},\ldots,W^{(k)}$. A
\emph{correlational multi-factor model} for an observable
$X=(X^{(1)},\ldots,X^{(d)})$ is a system of It\^o SDEs
\begin{equation}\label{eq:corr-SDE}
  dX^{(i)}_{t}\;=\;\mu_{i}(t,X_{t})\,dt
   \;+\;\sum_{j=1}^{k}\sigma_{ij}(t,X_{t})\,dW^{(j)}_{t},
   \qquad i=1,\ldots,d,
\end{equation}
with drift coefficients $\mu_{i}$ and diffusion coefficients
$\sigma_{ij}$. The instantaneous covariance is
$d\langle X^{(i)},X^{(j)}\rangle_{t}
 =\sum_{k}\sigma_{ik}\sigma_{jk}\,dt$, and the correlation
matrix $(\rho_{ij})$ is the central structural input for
co-movement. Hedging proceeds via delta-replication
along the joint Brownian filtration, with the classical Wiener
Clark--Ocone formula supplying the predictable
integrands~\cite[\S 7.1]{OperatorClarkOcone}.

A \emph{compositional model} for $X$ is one in which $X$ is the
value of a random field $\Xi$ sampled at a random argument
$\tau$:
\begin{equation}\label{eq:comp-X}
  X_{t}\;=\;\Phi\bigl(\Xi_{\tau_{t}}\bigr),
\end{equation}
where $\Xi=(\Xi_{u})_{u\in U}$ is a random measurable map from
a \emph{state space} $U$ to $\R$, $\tau=(\tau_{t})_{t\ge 0}$ is
a random $U$-valued process (the clock/state/sampling process),
$\Xi$ and $\tau$ are independent (or carry a joint dependence
structure separately specified), and $\Phi$ is a deterministic
post-processing.

\begin{example}[Random-field-sampled-by-process]\label{ex:field-sample}
Most generic compositional models have the form
\eqref{eq:comp-X} with $\Phi$ the identity:
$X_{t}=\Xi_{\tau_{t}}$. The map $u\mapsto\Xi_{u}$ is a fixed
sample-realized random landscape; $t\mapsto\tau_{t}$ is a
sampling trajectory in $U$; $X$ reads the landscape along the
trajectory.
\end{example}

The correlational model glues drivers linearly via the diffusion
coefficients $\sigma_{ij}$ and the correlation matrix
$\rho_{ij}$. Composition is structurally different in three
ways.

\begin{itemize}
\item \emph{Filtration enlargement.} The natural filtration of
the compositional observable $X_{t}=\Xi_{\tau_{t}}$ is
$\cF^{X}_{t}=\sigma(X_{s}:s\le t)$, which is strictly contained
in the source filtration $\sigma(\Xi)\vee\sigma(\tau)$.
Information about $\Xi$ outside the trajectory $\tau_{[0,t]}$,
and information about $\tau$ beyond the values reached, is
filtered out. The classical correlational filtration
$\sigma(W^{(1)},\ldots,W^{(k)}\restriction[0,t])$ has no such
collapse.

\item \emph{Topology.} The smoothness of $X$ in $t$ inherits
the roughness of both $\Xi$ along the trajectory $\tau$ and
$\tau$ itself. Even when $\Xi$ and $\tau$ are individually
Gaussian and $C^{0}$, the composition $\Xi_{\tau_{t}}$ can fail
to be a semimartingale (IBM is the canonical
example~\cite{EnergyFirst}); standard It\^o calculus does not
apply directly.

\item \emph{Hedge sensitivities.} A perturbation of $X$ at time
$t$ can come from a perturbation of $\Xi$ in a neighborhood of
$\tau_{t}$, or from a perturbation of $\tau$ that moves the
sampling argument. These two sensitivities have different
analytic content (the first is a local perturbation of the
field, the second involves the gradient of the field at the
random sampling point) and call for distinct hedge instruments.
\end{itemize}

In the correlational model these distinctions collapse: there
is one filtration, one quadratic-variation topology, one
delta-replication structure. In a compositional model they are
intrinsic.

%==============================================================
\section{IBM prototype}\label{sec:IBM-prototype}
%==============================================================

Let $\Bvar=(\Bvar_{x})_{x\in\R}$ be two-sided standard Brownian
motion in the spatial variable, and let
$\Yclock=(\Yclock_{s})_{s\ge 0}$ be an independent standard
Brownian motion. The product Gaussian space
\[
  (\Omega,\cF,\Prob)\;=\;(\Omega_{\Bvar}\times\Omega_{\Yclock},
   \cF_{\Bvar}\otimes\cF_{\Yclock},
   \Prob_{\Bvar}\otimes\Prob_{\Yclock})
\]
carries the iterated Brownian motion
\begin{equation}\label{eq:IBM-def}
  \IBM_{t}\;:=\;\Bvar(\Yclock_{t}),\qquad t\ge 0.
\end{equation}
$\Bvar$ is the random field (the \emph{scenery}); $\Yclock$ is
the random clock/state at which the field is sampled.
Cylindrical IBM observables of the form
\begin{equation}\label{eq:IBM-payoff}
  F\;=\;f(\Bvar(\Yclock_{T}))=f(\IBM_{T}),
   \qquad f\in C_{b}^{\infty}(\R),\;T>0,
\end{equation}
are the simplest payoff functionals of the composite observable
at a fixed maturity $T$. The analytic derivations of
\cref{sec:operator-derivatives,sec:carre,sec:operator-hedging}
hold under $f\in C_{b}^{\infty}(\R)$. The numerical experiment
of \cref{sec:experiment} additionally includes polynomial
payoffs (quadratic) as a stress test under the weaker
hypothesis that $f$ is smooth with sufficient integrability
against the Gaussian laws of $X$ and its perturbations for the
displayed identities to hold; bounded support is not required
for the empirical comparison.

IBM exposes the compositional architecture in the cleanest possible
setting. Both factors $\Bvar$ and $\Yclock$ are Gaussian and canonical,
with no nonlinearity in the SDE, no leverage, no calibration, and no
microstructure noise. The product source space admits a classical
Wiener representation, so one may also describe it using the usual
white-noise Malliavin calculus~\cite{NualartMalliavin}. That
representation is not the primitive calculus used here. The derivatives
in \cref{sec:DB-DY} are field/clock operator derivatives on the
compositional source coordinates. As
\cref{rem:operator-vs-malliavin} explains, the classical white-noise
Malliavin derivative is obtained only after pulling these coordinate
derivatives through the Gaussian covariance/kernel representation, and
the two should not be identified. The compositional pathology is
fully visible: $\IBM$ is not a semimartingale in its own
filtration~\cite{EnergyFirst}; the natural It\^o correction
$C_{\IBM,\varepsilon}(f;t)$ has $L^{2}$-divergent first
scenery-chaos as
$\varepsilon\to 0$~\cite[Prop.~5.1]{EnergyFirst}, closed in the
negative-Sobolev / quartic dual by~\cite{FCDC}; the joint
Wiener bidegree atlas is computed by~\cite{BidegreeAtlas}. IBM
is a mathematical prototype, not a model of asset prices.

%==============================================================
\section{Operator-derivative decomposition}\label{sec:operator-derivatives}
%==============================================================

\subsection{Field and clock derivatives}\label{sec:DB-DY}

For $F\in\mathrm{Dom}(D^{\mathrm{IBM}})$ on the product source space,
the \emph{source operator derivative}
of~\cite[\S 7.2]{OperatorClarkOcone} is the pair
\begin{equation}\label{eq:DIBM}
  D^{\mathrm{IBM}}F\;=\;(D^{\Bvar}F,D^{\Yclock}F),
\end{equation}
where $D^{\Bvar}_{x}F$ is the field/scenery operator derivative
at coordinate $x\in\R$ and $D^{\Yclock}_{s}F$ is the
clock/state operator derivative at time $s\in[0,T]$. The
field/scenery derivative $D^{\Bvar}_{x}F$ is the variational
derivative of the payoff with respect to perturbations of the
realized scenery value $\Bvar(x)$ at coordinate $x$; the
clock/state derivative $D^{\Yclock}_{s}F$ is the variational
derivative with respect to perturbations of the clock path
$\Yclock(s)$ at time $s$.

\begin{remark}[Operator derivative versus classical Malliavin
derivative]\label{rem:operator-vs-malliavin}
The derivative $D^{\Bvar}$ used here is the scenery-coordinate
operator derivative: it differentiates the payoff with respect
to perturbations of the realized field value $\Bvar(x)$ at
coordinate $x$, yielding, after mollification, the delta-like
kernel $\rho_{\varepsilon}(\Yclock_{T}-x)$
of~\eqref{eq:DB-eps} below. It is \emph{not} the classical
Malliavin derivative with respect to the white-noise increments
generating the Brownian scenery. For $\Bvar_{y}=W(\mathbf
1_{[0,y]})$ with $W$ scalar white noise on $\R$, the classical
Malliavin derivative is the interval kernel $D^{W}_{u}\Bvar_{y}
=\mathbf 1_{[0,y]}(u)$, structurally distinct from a delta. The
two are related by the covariance/kernel map of the underlying
Gaussian representation: in special cylindrical situations the
classical Malliavin derivative is the pullback of the
scenery-coordinate operator derivative through this map, but
the two live in different coordinate systems and should not be
identified. The precise relationship is characterized
in~\cite{OperatorClarkOcone}; the operator derivative is the
more primitive object and is what the compositional features
below are built from. The clock derivative $D^{\Yclock}$ admits the analogous variational
definition. For the cylindrical payoffs considered here, the same
kernel formula is obtained after representing the clock coordinate
through its Gaussian white-noise kernel. As with the scenery
coordinate, this is a representation-level pullback, not an
identification of the operator derivative with the classical
Malliavin derivative.
\end{remark}

For the cylindrical payoff $F=f(\Bvar(\Yclock_{T}))$ of
\eqref{eq:IBM-payoff}, the operator derivatives have explicit
expressions. \emph{Formally:}
\begin{align}
  D^{\Bvar}_{x}F
  &\;=\;f'\bigl(\Bvar(\Yclock_{T})\bigr)\,\delta_{\Yclock_{T}}(x)
   \quad\text{(distribution in $x$)},
   \label{eq:DB-formal}\\
  D^{\Yclock}_{s}F
  &\;=\;f'\bigl(\Bvar(\Yclock_{T})\bigr)\,\Bvar'(\Yclock_{T})\,\mathbf 1_{[0,T]}(s)
   \quad\text{(formal: $\Bvar'$ is white noise)}.
   \label{eq:DY-formal}
\end{align}
The scenery derivative is a delta-function in space at the
random point $\Yclock_{T}$ (a local-time/distributional
object); the clock derivative involves the spatial derivative
of Brownian scenery, which is white noise and hence not an
ordinary random variable.

\subsection{Mollification and the rough asymmetry}\label{sec:mollified}

Fix a mollifier
$\rho_{\varepsilon}(x)=\varepsilon^{-1}\rho(x/\varepsilon)$
(even, nonnegative, $C_{c}^{\infty}$, $\int\rho=1$,
$\supp\rho\subset[-1,1]$), and set
$\Bvar_{\varepsilon}:=\rho_{\varepsilon}*\Bvar$, so that
$\Bvar_{\varepsilon}\in C^{\infty}(\R)$ a.s. The mollified
cylindrical observable is
\begin{equation}\label{eq:F-eps}
  F_{\varepsilon}\;:=\;f\bigl(\Bvar_{\varepsilon}(\Yclock_{T})\bigr)\;\in\;\mathrm{Dom}(D^{\mathrm{IBM}}).
\end{equation}
The rigorous operator derivatives are
\begin{align}
  D^{\Bvar}_{x}F_{\varepsilon}
  &\;=\;f'\bigl(\Bvar_{\varepsilon}(\Yclock_{T})\bigr)\,\rho_{\varepsilon}(\Yclock_{T}-x),
   \label{eq:DB-eps}\\
  D^{\Yclock}_{s}F_{\varepsilon}
  &\;=\;f'\bigl(\Bvar_{\varepsilon}(\Yclock_{T})\bigr)\,\partial_{x}\Bvar_{\varepsilon}(\Yclock_{T})\,\mathbf 1_{[0,T]}(s).
   \label{eq:DY-eps}
\end{align}
Both are well-defined elements of $L^{2}(\Omega\times\R)$ and
$L^{2}(\Omega\times[0,T])$ respectively, for every
$\varepsilon>0$.

\Cref{eq:DB-eps,eq:DY-eps} make the two risk directions
explicit. The field-derivative
$D^{\Bvar}_{x}F_{\varepsilon}=f'(\Bvar_{\varepsilon}(\Yclock_{T}))\rho_{\varepsilon}(\Yclock_{T}-x)$
is concentrated near the random point $\Yclock_{T}$ and
measures field sensitivity in a neighborhood of size
$\varepsilon$ around the sampled spatial location. As
$\varepsilon\downarrow 0$, the spatial-mollified
$\rho_{\varepsilon}(\Yclock_{T}-x)$ converges to
$\delta_{\Yclock_{T}}(x)$ in the distributional sense in $x$;
the field derivative lives in a Sobolev space of negative
spatial order.

The clock-derivative
$D^{\Yclock}_{s}F_{\varepsilon}=f'(\Bvar_{\varepsilon}(\Yclock_{T}))\partial_{x}\Bvar_{\varepsilon}(\Yclock_{T})\mathbf 1_{[0,T]}(s)$
carries the more delicate object
$\partial_{x}\Bvar_{\varepsilon}(\Yclock_{T})$, the spatial
derivative of the smoothed scenery evaluated at the clock. As
$\varepsilon\downarrow 0$, the conditional variance scales as
\begin{equation}\label{eq:clock-var-scaling}
  \Var\bigl(\partial_{x}\Bvar_{\varepsilon}(\Yclock_{T})\bigm|\Yclock\bigr)
   \;=\;\varepsilon^{-1}\,\|\rho\|^{2}_{L^{2}}
\end{equation}
by the standard short-correlation scaling of mollified white
noise. This is the precise location of the IBM operator-energy
topology issue identified in~\cite[\S 5--6]{EnergyFirst}: the
unmollified clock-derivative integrand does not exist in
$L^{2}(\Omega)$ but only in the dual of a strictly weaker
operator-energy module, the quartic observable-energy dual
analyzed in~\cite{EnergyFirst,FCDC,BidegreeAtlas}.

In summary: $D^{\Bvar}_{x}F_{\varepsilon}$ is finite in
classical $L^{2}$ for every $\varepsilon>0$, with unmollified
limit in negative spatial Sobolev (a familiar
local-time/Bouleau--Yor type object);
$D^{\Yclock}_{s}F_{\varepsilon}$ is finite in classical $L^{2}$
for every $\varepsilon>0$, with unmollified limit in a strictly
weaker operator-energy module
(\cite{EnergyFirst,FCDC,BidegreeAtlas}). The asymmetry between
field and clock derivatives is the mathematical signature of
rough compositional risk: the sampling-process derivative
carries the field's spatial roughness through the chain rule.

%==============================================================
\section{Compositional carr\'e-du-champ}\label{sec:carre}
%==============================================================

\subsection{Product-source energy}\label{sec:carre-energy}

The operator-derivative framework equips the product source
space with a field/clock carr\'e-du-champ form built from the
operator derivatives $(D^{\Bvar},D^{\Yclock})$ of
\cref{sec:DB-DY},
\begin{equation}\label{eq:Gamma-IBM}
  \Gamma^{\mathrm{IBM}}(F,F)
   \;=\;\int_{\R}|D^{\Bvar}_{x}F|^{2}\,dx
       \;+\;\int_{0}^{T}|D^{\Yclock}_{s}F|^{2}\,ds,
   \qquad F\in\mathrm{Dom}(D^{\mathrm{IBM}}),
\end{equation}
with energy
\begin{equation}\label{eq:E-IBM}
  \cE^{\mathrm{IBM}}(F,F)\;=\;\E\bigl[\Gamma^{\mathrm{IBM}}(F,F)\bigr]
   \;=\;\E\!\int_{\R}|D^{\Bvar}_{x}F|^{2}\,dx
        \;+\;\E\!\int_{0}^{T}|D^{\Yclock}_{s}F|^{2}\,ds.
\end{equation}
Both integrals are finite for $F\in\mathrm{Dom}(D^{\mathrm{IBM}})$,
and $\Gamma^{\mathrm{IBM}}$ is the carr\'e-du-champ of the
operator-derivative pair $(D^{\mathrm{IBM}},\delta^{\mathrm{IBM}})$
via~\cite[Eq.~(7.5)]{OperatorClarkOcone}. The two terms
correspond to the two risk directions: field energy and
clock/state energy.

\subsection{Mollified vs unmollified}\label{sec:carre-moll}

For mollified payoffs $F_{\varepsilon}$ of~\eqref{eq:F-eps},
$\cE^{\mathrm{IBM}}(F_{\varepsilon},F_{\varepsilon})$ is finite
at every $\varepsilon>0$. As $\varepsilon\downarrow 0$, both
the field- and clock-energy integrals diverge like
$\varepsilon^{-1}$. For the field term,
\begin{equation}\label{eq:field-energy-blowup}
  \E\!\int_{\R}|D^{\Bvar}_{x}F_{\varepsilon}|^{2}\,dx
   \;=\;\E\bigl[(f'(\Bvar_{\varepsilon}(\Yclock_{T})))^{2}\bigr]
        \cdot\int_{\R}\rho_{\varepsilon}(\Yclock_{T}-x)^{2}\,dx
   \;=\;\E\bigl[(f'(\Bvar_{\varepsilon}(\Yclock_{T})))^{2}\bigr]
        \cdot\varepsilon^{-1}\,\|\rho\|_{L^{2}}^{2},
\end{equation}
which diverges as $\varepsilon^{-1}$ provided $f'$ does not
vanish identically on the support of
$\Bvar_{\varepsilon}(\Yclock_{T})$. For the clock term,
\begin{equation}\label{eq:clock-energy-blowup}
  \E\!\int_{0}^{T}|D^{\Yclock}_{s}F_{\varepsilon}|^{2}\,ds
   \;=\;T\cdot\E\bigl[(f'(\Bvar_{\varepsilon}(\Yclock_{T})))^{2}\,(\partial_{x}\Bvar_{\varepsilon}(\Yclock_{T}))^{2}\bigr]
   \;\lesssim\;T\cdot\varepsilon^{-1}\,\|f'\|_{L^{\infty}}^{2}\,\|\rho\|_{L^{2}}^{2},
\end{equation}
also diverging as $\varepsilon^{-1}$ by the same
short-correlation scaling, with a matching lower bound under the
nondegeneracy condition $\E[(f'(\Bvar_{\varepsilon}(\Yclock_{T})))^{2}]
\gtrsim 1$. \Cref{eq:clock-energy-blowup}
is~\eqref{eq:clock-var-scaling} integrated against the clock-time
variable.

The field/clock asymmetry is therefore not finiteness of the
$L^{2}$ energies---both diverge---but the \emph{nature of the
distributional limits}. The field derivative
$D^{\Bvar}_{x}F_{\varepsilon}=f'(\Bvar_{\varepsilon}(\Yclock_{T}))\rho_{\varepsilon}(\Yclock_{T}-x)$
converges, distributionally in $x$, to $f'(\Bvar(\Yclock_{T}))\,\delta_{\Yclock_{T}}(x)$:
a clean local-time/Bouleau--Yor-type object in negative spatial
Sobolev. The clock derivative carries the additional rough
factor $\partial_{x}\Bvar_{\varepsilon}(\Yclock_{T})$ through
the chain rule, which does not admit a comparable clean
distributional limit; it lands instead in the strictly weaker
quartic operator-energy module
of~\cite{EnergyFirst,FCDC,BidegreeAtlas}.

The compositional carr\'e-du-champ $\Gamma^{\mathrm{IBM}}$ thus
does not converge to a classical $L^{2}$-form as
$\varepsilon\downarrow 0$. The natural target is the
negative-Sobolev / quartic operator-energy dual identified
in~\cite{EnergyFirst,FCDC,BidegreeAtlas}: the unmollified
clock-derivative integrand and its dual integrals are naturally
treated in this negative-Sobolev / quartic operator-energy dual
rather than in $L^{2}(\Omega)$. In compositional-risk language, field-spatial
roughness, transported by chain rule through the random clock,
requires a renormalized topology to make sense of clock-state
sensitivities at the unmollified limit.

\begin{remark}[Rough sampling, smooth field, would be the
opposite case]\label{rem:rough-sampling-smooth-field}
The asymmetry between the field-energy and clock-energy terms
in~\eqref{eq:E-IBM} reflects that, in the IBM
$\IBM_{t}=\Bvar(\Yclock_{t})$, the field $\Bvar$ is rough
($\partial_{x}\Bvar$ is white noise). Were the field smooth and
the clock rough, the asymmetry would reverse. In a general
compositional model $X_{t}=\Phi(\Xi_{\tau_{t}})$, both kinds of
roughness can interact and the natural carr\'e-du-champ is
correspondingly more complex; IBM is the cleanest case where
the analysis is tractable end-to-end.
\end{remark}

%==============================================================
\section{Operator hedging: motivation for the feature set}\label{sec:operator-hedging}
%==============================================================

The operator Clark--Ocone formula
of~\cite[\S 7]{OperatorClarkOcone} applied to a mollified IBM
payoff $F_{\varepsilon}=f(\Bvar_{\varepsilon}(\Yclock_{T}))$
gives, in compact form,
\begin{equation}\label{eq:hedge-IBM-source}
  F_{\varepsilon}-\E F_{\varepsilon}
   \;=\;\delta_{\Bvar}\bigl(\Piad_{\Bvar}D^{\Bvar}F_{\varepsilon}\bigr)
        \;+\;\delta_{\Yclock}\bigl(\Piad_{\Yclock}D^{\Yclock}F_{\varepsilon}\bigr),
\end{equation}
where $\delta_{\Bvar},\delta_{\Yclock}$ are the Skorohod
divergences against the scenery and clock Brownian factors and
$\Piad_{\Bvar},\Piad_{\Yclock}$ are the predictable projections
in the appropriate enlarged filtrations
(\cite[\S 7.3]{OperatorClarkOcone}). The point extracted from
this formula is not a tradable hedging strategy; it tells us
which features the centered payoff is a function of.

Reading the chain rule on the cylindrical IBM payoff
(\cref{eq:F-eps,eq:DB-eps,eq:DY-eps}), the integrands in
\eqref{eq:hedge-IBM-source} are linear combinations of six
state-conditioned compositional features:
\begin{equation}\label{eq:op-feature-set}
  \delta\Bvar_{\varepsilon}(\Yclock),\quad
  \partial_{x}\Bvar_{\varepsilon}(\Yclock)\,\delta\Yclock,\quad
  f'(X)\,\delta\Bvar_{\varepsilon}(\Yclock),\quad
  f'(X)\,\partial_{x}\Bvar_{\varepsilon}(\Yclock)\,\delta\Yclock,\quad
  \Gamma_{\mathrm{field}},\;\Gamma_{\mathrm{clock}},
\end{equation}
where
$\Gamma_{\mathrm{field}}=(f'(X)\,\delta\Bvar_{\varepsilon}(\Yclock))^{2}$
and
$\Gamma_{\mathrm{clock}}=(f'(X)\,\partial_{x}\Bvar_{\varepsilon}(\Yclock)\,\delta\Yclock)^{2}$
are the squared compositional carr\'e-du-champ terms
(\cref{sec:carre}). The first four are the analytic integrands
of the operator Clark--Ocone formula; the last two are the
natural second-order energy corrections.

The features \eqref{eq:op-feature-set} are the central object
of the empirical section. \Cref{sec:experiment} compares a
fixed ML model trained on this feature set (plus a generic
non-compositional baseline) against the same model trained on
the baseline alone, on compositional and on correlation-matched
data.

\begin{remark}[Unmollified limit and natural IBM filtration]\label{rem:hedge-not-the-point}
Equation~\eqref{eq:hedge-IBM-source} lives on the product
source filtration $\sigma(\Bvar)\vee\sigma(\Yclock)$;
restriction to the natural IBM observable filtration
$\sigma(\IBM_{s}:s\le t)$ is a separate non-semimartingale
problem~\cite{EnergyFirst} not addressed here. The unmollified
limit of the clock integrand exists only in an operator-energy
module strictly weaker than $L^{2}$, identified
in~\cite{EnergyFirst,FCDC,BidegreeAtlas}. The empirical content
concerns feature engineering for ML on the mollified payoff at
fixed $\varepsilon$, not analytic hedging at the unmollified
limit.
\end{remark}

%==============================================================
\section{Finance interpretation}\label{sec:finance}
%==============================================================

The compositional model \eqref{eq:comp-X} captures several
market-risk situations in which the standard correlational
formulation~\eqref{eq:corr-model} is structurally inadequate.
The unifying finance punchline is the field/clock split:
\emph{field risk} is sensitivity to the random market surface
$\Xi$ at the current sampling location; \emph{clock/state risk}
is sensitivity to moving through that surface, i.e.\ to
perturbations of the sampling-state $\tau$; correlation-only
models collapse the two into a single off-diagonal covariance
entry and lose the distinction. Four illustrative classes,
qualitative not calibrated:

\begin{itemize}
\item \emph{Latent value field sampled by order flow.}
$\Xi_{u}$ is a latent fundamental-value or risk-premium surface
indexed by a market state-space $u\in U$; $\tau_{t}$ is the
order-flow process that determines which state is expressed in
observed prices. Field risk is shifts in the latent value
surface; clock/state risk is order flow sampling a different
region of the surface even when the surface is unchanged. A
correlation-only model conflates these into a single
quadratic-variation contribution.
\item \emph{Volatility surface sampled by liquidity regime.}
$\Xi_{u}$ is an implied- or local-volatility surface as a
function of strike, maturity, or some abstract liquidity
parameter $u$; $\tau_{t}$ is a random regime-state process.
Volatility risk splits into surface-shape risk
(field/scenery) and regime-shift risk (clock/state). Hedging a
volatility option requires distinct instruments for each.
\item \emph{Execution-cost field sampled by trade intensity.}
$\Xi_{u}$ is an instantaneous execution-cost map (as a function
of inventory, urgency, or order-book depth); $\tau_{t}$ is the
cumulative trading-intensity process. The realized
execution-cost stream depends on field shifts (changes in
market depth and impact) and on clock shifts (changes in
trade-intensity behavior). Optimal-execution in this language
separates field-risk hedging (e.g.\ via market-maker rebates)
from clock-risk hedging (e.g.\ via algorithmic re-pacing).
\item \emph{Credit / liquidity state sampled by event clock.}
$\Xi_{u}$ is a credit-/liquidity-state landscape (a matrix of
recovery rates or default-intensity values) indexed by an
abstract risk-state; $\tau_{t}$ is a random-event clock
triggering transitions between states. Field risk is variation
of recovery/intensity at fixed event clock; clock/state risk is
variation of the event clock at fixed field. Distress events
typically involve both.
\end{itemize}

In each case, a covariance-based decomposition collapses field
risk and clock/state risk into a single off-diagonal covariance
entry; the operator-derivative decomposition
$D=(D^{\Xi},D^{\tau})$ separates them along their distinct
risk directions. The compositional carr\'e-du-champ
of \cref{sec:carre} provides the corresponding energy
decomposition; the operator hedging formula
of \cref{sec:operator-hedging} provides the corresponding
hedge decomposition.

\begin{remark}[Why this is not stochastic volatility, time
change, or regime-switching]\label{rem:not-just}
Each of the four examples could be locally approximated by a
correlational stochastic-volatility, time-change, or
regime-switching specification. The compositional formulation
is structurally distinct in three ways flagged
in~\cref{sec:correl-vs-comp}: filtration enlargement,
non-semimartingale topology, and the field/clock-hedge
asymmetry made precise by the operator derivatives. These
distinctions are not visible to covariance/PCA features unless
the compositional architecture is explicitly modeled; they are
the mathematical content of treating a payoff as a composition
rather than as a smooth transform of correlated drivers.
\end{remark}

%==============================================================
\section{Numerical experiment: compositional features for ML}\label{sec:experiment}
%==============================================================

This experiment tests the framework directly. The ML model and
training protocol are held fixed; the data and the train/test
split are held fixed; the only variable is the feature set. The
operator-derivative framework and ML are complementary, not
competing: operator derivatives identify compositional features;
ML learns nonlinear combinations and higher-order corrections.

Two questions are tested separately. (i) On compositional data
in the mid-regime ($\eta_{Y}=\eta_{B}=0.1$), can ML on the
operator-derived features
\eqref{eq:op-feature-set} match the analytic first-order
operator hedge, and outperform ML on non-compositional features
(PCA of field shocks, $X$, $f'(X)$, $Y$, plus standard
interactions)? The answer is near-definitional, since the
operator features \emph{are} the first-order Taylor terms
of $\Delta F$ in the source shocks: in the mid-regime ML on
operator features reaches $R^{2}\approx 0.96$, essentially
reproducing the analytic operator hedge. (ii) In the
large-clock-shock regime where the first-order Taylor breaks
down, does ML on operator features extract higher-order
content beyond what the analytic first-order operator hedge
can deliver? This is the empirically informative question and
the centerpiece of the experiment: at $\eta_{Y}=0.4$, the
analytic first-order operator hedge collapses to
$R^{2}\approx 0.51$ while gradient-boosted ML on operator
features (including the squared $\Gamma$ terms) holds at
$0.65$, a $+0.14$ recovery that the first-order formula cannot
produce.

\paragraph{Setup.}
We simulate a mollified IBM payoff
$F_{\varepsilon}=f(\Bvar_{\varepsilon}(\Yclock_{T}))$
(\cref{sec:IBM-prototype}) on a spatial grid $[-8,8]$ with
$\Delta x=0.1$, $\Yclock\sim N(0,\sigma_{Y}^{2}T)$ with $T=1$
and $\sigma_{Y}=1$, Gaussian mollifier of bandwidth
$\varepsilon=0.3$, and a one-step P\&L shock
$\Delta F=f((\Bvar+\eta_{B}\delta\Bvar)_{\varepsilon}(\Yclock+\eta_{Y}\delta\Yclock))-F$
with $\eta_{B}=\eta_{Y}=0.1$. For the correlation-matched
negative control, $X^{\mathrm{corr}}=aZ_{1}+bZ_{2}$ with
$a^{2}=b^{2}=\Var(\Bvar_{\varepsilon}(\Yclock))/2$. Each seed
draws $n_{\mathrm{paths}}=10{,}000$ paths and splits
50/50 train/test. We run 12 seeds for the main payoff sweep.

This is not time-series forecasting; it is controlled
out-of-sample approximation of the P\&L shock map conditional
on source shocks $(\delta\Bvar,\delta\Yclock)$. The shocks are
inputs to the model, not predictions of the model; the
regression target is $\Delta F$ as a function of those source
shocks together with the pre-shock state.

\paragraph{Code and data.}
Simulation code, seeds, hyperparameters, and CSV outputs are
available at
\url{https://github.com/QuijoticResearch/compositional-risk-features}.
\Cref{tab:code-mapping} maps each table and figure below to the
driver script that produces it; intermediate CSVs are written to
\texttt{experiments/results/} and figures to
\texttt{experiments/figures/}.

\begin{table}[ht]
\centering
\footnotesize
\renewcommand{\arraystretch}{1.15}
\begin{tabularx}{\linewidth}{@{}>{\raggedright\arraybackslash}p{0.30\linewidth} >{\raggedright\arraybackslash}p{0.32\linewidth} >{\raggedright\arraybackslash}X@{}}
\toprule
Result & Driver script (\path{experiments/}) & Outputs (\path{experiments/results/}) \\
\midrule
\Cref{tab:main-results}, \Cref{tab:delta-feature}, \Cref{fig:main}
   & \path{compositional_hedge_sim_v3.py}
   & \path{v3_main_results.csv}, \path{v3_main_summary.csv} \\
\Cref{tab:negctrl}
   & \path{compositional_hedge_sim_v3_negctrl.py}
   & \path{v3_negctrl_*.csv} \\
\Cref{tab:nplus}
   & \path{compositional_hedge_sim_v4_nplus.py}
   & \path{v4_nplus_main_*.csv} \\
\Cref{tab:stress-etaY}
   & \path{compositional_hedge_sim_v4_nplus.py}
   & \path{v4_nplus_stress_etaY.csv} \\
\Cref{tab:stress-eps}
   & \path{compositional_hedge_sim_v3.py}
   & \path{v3_stress_eps.csv} \\
\Cref{fig:stress}
   & \path{compositional_hedge_sim_v3.py}
   & \path{v3_stress_etaY.csv}, \path{v3_stress_eps.csv} \\
\Cref{fig:importance}
   & \path{compositional_hedge_sim_v3.py}
   & \path{v3_feature_importance.csv} \\
\bottomrule
\end{tabularx}
\caption{Mapping of empirical results to driver scripts in the
companion repository. All scripts run on CPU; total compute for
the headline numbers is on the order of minutes.}
\label{tab:code-mapping}
\end{table}

\paragraph{Payoffs.}
The four payoffs are fixed throughout. Softplus and smoothed
call use $f(x)=\beta^{-1}\log(1+e^{\beta(x-K)})$ with
$(\beta,K)=(2,0)$ and $(4,\tfrac12)$ respectively (the call
parameters give a Black--Scholes-like elbow at the money);
quadratic is $f(x)=x^{2}$ and tanh is $f(x)=\tanh(x)$.
\Cref{tab:payoffs} collects the formulas and parameters.

\begin{table}[ht]
\centering
\begin{tabular}{lll}
\toprule
Payoff & $f(x)$ & Parameters \\
\midrule
softplus       & $\beta^{-1}\log(1+e^{\beta(x-K)})$ & $\beta=2,\;K=0$ \\
smoothed call  & $\beta^{-1}\log(1+e^{\beta(x-K)})$ & $\beta=4,\;K=\tfrac12$ \\
quadratic      & $x^{2}$                            & --- \\
tanh           & $\tanh(x)$                         & --- \\
\bottomrule
\end{tabular}
\caption{Payoff functions used throughout
\cref{sec:experiment}. Softplus and tanh are bounded-derivative
smooth payoffs; smoothed call has the same softplus form with a
sharper $\beta$ and strike $K=\tfrac12$; quadratic is a
polynomial stress test with unbounded $f'$.}
\label{tab:payoffs}
\end{table}

\paragraph{Three method classes.}
\begin{itemize}
\item \emph{A. Pure operator hedge} (analytic, no fitting):
\[
  \widehat F^{\mathrm{op}}
   \;=\;f'(X)\bigl[\eta_{B}\,\delta\Bvar_{\varepsilon}(\Yclock)
        +\eta_{Y}\,\partial_{x}\Bvar_{\varepsilon}(\Yclock)\,\delta\Yclock\bigr].
\]
\item \emph{B. ML on non-compositional features} (set N): PCA
of $\delta\Bvar$ (top $K=10$),
$\delta\Bvar_{\varepsilon}(0)$, $\delta\Yclock$, $X$, $f'(X)$,
$\Yclock$, plus interactions
$f'(X)\delta\Yclock$, $f'(X)\delta\Bvar_{\varepsilon}(0)$,
$X\delta\Yclock$, $\Yclock\delta\Yclock$.
\item \emph{C. ML on operator features} (set Op): set N
augmented with the six state-conditioned compositional features
\eqref{eq:op-feature-set}.
\end{itemize}
ML models in classes B and C: Ridge ($\alpha=1$) and
\texttt{HistGradientBoostingRegressor} (HGB, $200$ trees,
learning rate $0.05$, max depth $6$); see
Friedman~\cite{FriedmanGBM} for the gradient boosting algorithm.

\paragraph{Main diagnostic.}
\begin{equation}\label{eq:delta-feature}
  \Delta R^{2}_{\mathrm{feature}}\;:=\;R^{2}\bigl(\text{best ML on Op}\bigr)\;-\;R^{2}\bigl(\text{best ML on N}\bigr),
\end{equation}
plus the auxiliary
$R^{2}(\text{best ML on Op})-R^{2}(\text{pure operator})$,
which measures whether ML on operator features improves on the
analytic first-order operator hedge.

\paragraph{Main results.}
\Cref{tab:main-results} summarizes the per-payoff $R^{2}$
across the five headline methods (12 seeds, four payoffs).

\begin{table}[ht]
\centering
\small
\begin{tabular}{lcccc}
\toprule
Method & softplus & quadratic & tanh & smoothed call \\
\midrule
ML+N, Ridge                & $0.156$ & $0.147$ & $0.156$ & $0.148$ \\
ML+N, HGB                  & $0.412$ & $0.315$ & $0.342$ & $0.427$ \\
\textbf{ML+Op, Ridge}      & $\mathbf{0.966}$ & $\mathbf{0.963}$ & $\mathbf{0.957}$ & $\mathbf{0.957}$ \\
ML+Op, HGB                 & $0.953$ & $0.916$ & $0.953$ & $0.948$ \\
pure operator              & $0.965$ & $0.964$ & $0.954$ & $0.957$ \\
\bottomrule
\end{tabular}
\caption{Main comparison, mean out-of-sample $R^{2}$ across 12
seeds, $n=10{,}000$, $\sigma_{Y}=1$, $\eta_{Y}=\eta_{B}=0.1$,
$\varepsilon=0.3$. ML+N (HGB) captures $0.32$--$0.43$ of
variance; ML+Op (Ridge) captures $0.96+$, essentially matching
the analytic operator hedge.}
\label{tab:main-results}
\end{table}

\begin{table}[ht]
\centering
\begin{tabular}{lccc}
\toprule
Payoff & $\Delta R^{2}_{\mathrm{feature}}$ & $R^{2}_{\mathrm{op,pure}}$ &
  $R^{2}_{\mathrm{ML+Op}}-R^{2}_{\mathrm{op,pure}}$ \\
\midrule
softplus      & $0.554 \pm 0.005$ & $0.965$ & $+0.001 \pm 0.000$ \\
quadratic     & $0.648 \pm 0.010$ & $0.964$ & $-0.001 \pm 0.001$ \\
tanh          & $0.616 \pm 0.005$ & $0.954$ & $+0.003 \pm 0.000$ \\
smoothed call & $0.531 \pm 0.007$ & $0.957$ & $+0.001 \pm 0.001$ \\
\bottomrule
\end{tabular}
\caption{Main diagnostics, mean$\pm$SE across 12 seeds.
\emph{Left:} $\Delta R^{2}_{\mathrm{feature}}$ = (best ML on
Op) $-$ (best ML on N): supplying ML with operator features
adds $0.53$--$0.65$ units of $R^{2}$ across payoffs.
\emph{Right:} ML on operator features does not improve on the
pure analytic operator hedge in the mid-regime; it essentially
reproduces it ($\pm 0.003$).}
\label{tab:delta-feature}
\end{table}

\paragraph{Correlation-matched negative control.}
The operator features may simply be better-engineered features
that would help ML in any setting, not just on compositional
data. We falsify this by replacing the regression target with
a correlation-matched analogue. The features (set N and set
Op) are computed exactly as before from the same simulated
$(\Bvar,\Yclock)$ paths. The target $\Delta F$ is replaced by
$\Delta F^{\mathrm{corr}}=f(X^{\mathrm{corr}}+\Delta X^{\mathrm{corr}})-f(X^{\mathrm{corr}})$
with $X^{\mathrm{corr}}\sim N(0,\Var(X))$ and
$\Delta X^{\mathrm{corr}}\sim N(0,\Var(\Delta F)/\E[f'(X)^{2}])$
drawn independently of $(\Bvar,\Yclock)$, so $X^{\mathrm{corr}}$
has the same marginal variance as the compositional state but
carries no dependence on the field/clock structure. The
operator features should now carry no information about the
target. \Cref{tab:negctrl} confirms this prediction: the
compositional-data $\Delta R^{2}_{\mathrm{feature}}\approx 0.6$
collapses to $|\Delta R^{2}_{\mathrm{feature}}|\le 0.01$ on the
correlation-matched target, across all four payoffs. The
advantage is tied to compositional structure rather than
generic nonlinear feature enrichment under the tested HGB
protocol.

\begin{table}[ht]
\centering
\begin{tabular}{lcc}
\toprule
Payoff & $\Delta R^{2}_{\mathrm{feature}}$ (compositional) &
  $\Delta R^{2}_{\mathrm{feature}}$ (correlational) \\
\midrule
softplus      & $+0.609 \pm 0.010$ & $-0.003 \pm 0.003$ \\
quadratic     & $+0.634 \pm 0.012$ & $-0.002 \pm 0.002$ \\
tanh          & $+0.642 \pm 0.006$ & $+0.008 \pm 0.002$ \\
smoothed call & $+0.592 \pm 0.014$ & $-0.002 \pm 0.003$ \\
\bottomrule
\end{tabular}
\caption{Correlation-matched negative control,
$\Delta R^{2}_{\mathrm{feature}}=R^{2}(\text{HGB+Op})-R^{2}(\text{HGB+N})$,
mean$\pm$SE across 8 seeds, $n=6{,}000$. \emph{Left column:}
compositional target $\Delta F$, replicating the headline.
\emph{Right column:} target replaced by
$f(X^{\mathrm{corr}}+\Delta X^{\mathrm{corr}})-f(X^{\mathrm{corr}})$
with $X^{\mathrm{corr}}$ drawn independently of
$(\Bvar,\Yclock)$, marginal variance matched. The HGB+Op
advantage of $\sim 0.6$ on compositional data collapses to
noise level ($|\Delta|\le 0.01$) on correlation-matched data;
on correlational targets both ML+N and ML+Op overfit and reach
small negative $R^{2}$, so neither feature set explains
anything and there is nothing for the operator features to
add.}
\label{tab:negctrl}
\end{table}

\paragraph{State-conditioned baseline (Set N+).}
A separate objection is that the operator advantage may simply
reflect any state-conditioned feature set, not the
operator-derivative directions specifically. A strong baseline
that a practitioner without operator calculus could build is
\emph{locally} state-aware around the sampled clock position
$\Yclock$: PCA of $\delta\Bvar$ restricted to the 11 grid points
nearest $\Yclock$ (top six components), the five raw
nearby-grid shocks
$\delta\Bvar(\Yclock-2\Delta x),\ldots,\delta\Bvar(\Yclock+2\Delta x)$,
and twenty random Fourier features (RBFSampler with $\gamma=1$)
over $(X,\Yclock,\delta\Yclock,\delta\Bvar_{\varepsilon}(0),
\text{raw nearby shocks})$. Call this \emph{Set N+}: it
strictly extends Set N with state-conditioned, locally nonlinear
capacity, but contains \emph{no} operator-derivative directions
(no $f'(X)\,\delta\Bvar_{\varepsilon}(\Yclock)$, no clock
derivative, no $\Gamma$ squares). Set N+ has 50 features
against Set Op's 25; the comparison is whether more capacity in
the wrong directions can match fewer features in the right
directions. The prediction, recorded before running the N+
comparison, was that the gap shrinks against N+; the open
question was whether it shrinks to noise or holds at a
meaningful margin. \Cref{tab:nplus} shows that the gap shrinks
by roughly a third but remains $0.35$--$0.49$ in $R^{2}$ across
all four payoffs: the operator advantage is specifically about
\emph{which directions} in the source shock are exploited, not
merely \emph{where} in the field they are sampled.

\begin{table}[ht]
\centering
\small
\setlength{\tabcolsep}{5pt}
\begin{tabular}{lccccc}
\toprule
Payoff
 & \makecell{N\\(best ML)}
 & \makecell{N+\\(best ML)}
 & \makecell{Op\\(best ML)}
 & \makecell{Op $-$ N}
 & \makecell{Op $-$ N+} \\
\midrule
softplus       & $0.412$ & $0.557$ & $0.966$ & $+0.554 \pm 0.005$ & $+0.409 \pm 0.005$ \\
quadratic      & $0.315$ & $0.494$ & $0.963$ & $+0.648 \pm 0.010$ & $+0.469 \pm 0.007$ \\
tanh           & $0.342$ & $0.465$ & $0.957$ & $+0.616 \pm 0.005$ & $+0.492 \pm 0.006$ \\
smoothed call  & $0.427$ & $0.605$ & $0.957$ & $+0.531 \pm 0.007$ & $+0.353 \pm 0.007$ \\
\bottomrule
\end{tabular}
\caption{Three-arm comparison at the main regime
($\sigma_{Y}=1$, $\eta_{Y}=\eta_{B}=0.1$, $\varepsilon=0.3$,
12 seeds, $n=10{,}000$). ``Best ML'' is the maximum of Ridge
and HGB on the given feature set. Set N+ is strictly stronger
than Set N---it adds local PCA, raw nearby-grid shocks, and
RFF over local state---and closes
$0.12$--$0.18$ of the gap to Op. The residual Op $-$ N+ gap is
$0.35$--$0.49$, larger than a pure local-state-encoding account
of the operator features would predict.}
\label{tab:nplus}
\end{table}

\begin{figure}[ht]
\centering
\includegraphics[width=0.95\linewidth]{experiments/figures/v3_main_comparison.pdf}
\caption{Main comparison: mean out-of-sample $R^{2}$ on
$\Delta F$, per method and payoff. ML on non-compositional
features (left two groups) reaches $0.15$--$0.43$. ML on
operator features (middle two) and pure operator (rightmost)
all sit at $0.92$--$0.97$.}
\label{fig:main}
\end{figure}

\paragraph{Stress: large clock shock (the key result).}
The mid-regime headline of $R^{2}\approx 0.96$ on operator
features is, by construction, close to the maximum any
first-order shock-linear hedger could attain: at small
$(\eta_{B},\eta_{Y})$ the P\&L shock $\Delta F$ is a first-order
Taylor expansion in the source shocks, whose coefficients are
the operator features themselves. Reaching $R^{2}\approx 0.96$
in that regime confirms that the operator features expose the
first-order shock structure cleanly. The discriminating regime
is large clock shock, where the first-order operator hedge
breaks down and the question is whether ML on operator features
can recover the missing higher-order content.

\Cref{tab:stress-etaY} sweeps the clock-shock size $\eta_{Y}$
at fixed $\sigma_{Y}=1$, $\varepsilon=0.3$;
\cref{tab:stress-eps} sweeps $\varepsilon$ at fixed
$\sigma_{Y}=1$, $\eta_{Y}=0.1$. Eight seeds per cell,
$n_{\mathrm{paths}}=6{,}000$, softplus payoff. The decoupling
of historical clock variance $\sigma_{Y}$ from shock size
$\eta_{Y}$ matters: at $\eta_{Y}=0.4$ the analytic first-order
operator hedge collapses from $R^{2}\approx 0.97$ to $0.51$,
while HGB on operator features holds at $0.65$, a $+0.14$
recovery. This is the empirically informative cell: the
analytic hedge knows only the first-order direction, while
HGB, given the $\Gamma_{\mathrm{field}}$ and
$\Gamma_{\mathrm{clock}}$ second-order features, learns
higher-order corrections that the first-order formula cannot
produce.

\begin{table}[ht]
\centering
\small
\setlength{\tabcolsep}{4pt}
\begin{tabular}{lcccccc}
\toprule
$\eta_{Y}$
 & \makecell{ML+N\\HGB}
 & \makecell{ML+N+\\HGB}
 & \makecell{ML+Op\\HGB}
 & \makecell{ML+Op\\Ridge}
 & \makecell{pure\\op}
 & \makecell{ML+Op\\$-$ pure} \\
\midrule
$0.02$ & $0.583$ & $0.937$ & $0.983$ & $0.998$ & $0.997$ & $+0.001$ \\
$0.05$ & $0.480$ & $0.765$ & $0.970$ & $0.992$ & $0.992$ & $+0.000$ \\
$0.10$ & $0.338$ & $0.521$ & $0.946$ & $0.967$ & $0.966$ & $+0.001$ \\
$0.20$ & $0.264$ & $0.331$ & $0.858$ & $0.863$ & $0.851$ & $+0.012$ \\
$0.40$ & $0.242$ & $0.256$ & $\mathbf{0.649}$ & $0.617$ & $0.510$ & $\mathbf{+0.139}$ \\
\bottomrule
\end{tabular}
\caption{Stress over $\eta_{Y}$, with the state-conditioned N+
baseline included (8 seeds per cell, $n=6{,}000$, softplus
payoff). At very small $\eta_{Y}=0.02$ (deep Taylor regime), N+
catches up to within $0.06$ of Op because the local nonlinear
baseline can approximate the first-order shock structure that
the operator features encode. From $\eta_{Y}=0.05$ onward Op
opens a $0.22$--$0.53$ margin over N+; at $\eta_{Y}=0.4$ pure
operator collapses to $0.51$, while ML+Op (HGB) recovers $0.65$,
an improvement of $+0.14$ over the analytic first-order hedge,
and ML+N+ (HGB) holds only $0.26$.}
\label{tab:stress-etaY}
\end{table}

\begin{table}[ht]
\centering
\small
\setlength{\tabcolsep}{5pt}
\begin{tabular}{lccccc}
\toprule
$\varepsilon$
 & \makecell{ML+N\\HGB}
 & \makecell{ML+Op\\HGB}
 & \makecell{ML+Op\\Ridge}
 & \makecell{pure\\op}
 & \makecell{ML+Op\\$-$ pure} \\
\midrule
$0.15$ & $0.208$ & $0.870$ & $0.887$ & $0.877$ & $+0.010$ \\
$0.30$ & $0.338$ & $0.946$ & $0.967$ & $0.966$ & $+0.001$ \\
$0.60$ & $0.496$ & $0.964$ & $0.988$ & $0.988$ & $+0.000$ \\
$1.20$ & $0.634$ & $0.970$ & $0.994$ & $0.994$ & $+0.000$ \\
$2.40$ & $\mathbf{0.841}$ & $0.974$ & $0.971$ & $0.958$ & $+0.013$ \\
\bottomrule
\end{tabular}
\caption{Stress over $\varepsilon$ on a log axis. ML+Op (Ridge
or HGB) tracks or slightly exceeds the pure operator hedge
across the full $\varepsilon$ range. ML+N (HGB) catches up only
in the very smooth regime $\varepsilon=2.4$ where it reaches
$0.84$. In the rough regime $\varepsilon=0.15$ the structural
advantage $\mathrm{(ML+Op)}-\mathrm{(ML+N)}\approx 0.68$ is
largest; ML+Op even slightly improves on pure operator there.}
\label{tab:stress-eps}
\end{table}

\begin{figure}[ht]
\centering
\includegraphics[width=0.95\linewidth]{experiments/figures/v3_stress_curves.pdf}
\caption{Stress curves. Left: $R^{2}$ vs $\eta_{Y}$ at fixed
$\sigma_{Y}=1$, $\varepsilon=0.3$. Right: $R^{2}$ vs
$\varepsilon$ (log) at fixed $\sigma_{Y}=1$, $\eta_{Y}=0.1$.
ML+Op (HGB) overtakes the pure operator hedge at large
$\eta_{Y}$ by capturing higher-order corrections through the
$\Gamma$-square features; ML+N stays well below ML+Op except
in the very smooth-field regime.}
\label{fig:stress}
\end{figure}

\paragraph{Feature ablation.}
\Cref{fig:importance} reports the leave-one-out drop in $R^{2}$
when one operator feature is removed from the Op set (HGB
refit on the remaining features), averaged over 5 seeds per
payoff at the main-regime settings.

\begin{figure}[ht]
\centering
\includegraphics[width=0.92\linewidth]{experiments/figures/v3_feature_importance.pdf}
\caption{Leave-one-out drop in $R^{2}$ for HGB trained on Op
features, by dropped operator feature and by payoff. For the
quadratic payoff (orange), the two state-conditioned hedge
integrands $f'(X)\delta\Bvar_{\varepsilon}(\Yclock)$ and
$f'(X)\partial_{x}\Bvar_{\varepsilon}(\Yclock)\delta\Yclock$
are individually critical (drop of $\sim 0.13$--$0.19$ each);
for softplus, tanh, and smoothed call, the individual drops
are modest ($\le 0.01$) because the operator features carry
correlated information that ML can partially recover from the
remaining features. The collective contribution of the
operator features is large (see \cref{tab:delta-feature});
individual features matter most for the quadratic payoff where
the nonlinearity structure $f'(x)=2x$ is exposed cleanly.}
\label{fig:importance}
\end{figure}

\paragraph{Reading.}
\begin{itemize}
\item In the mid-regime ($\eta_{Y}=0.1$, $\varepsilon=0.3$),
ML on operator features matches the pure analytic operator
hedge to within sampling error. ML offers no improvement on
the first-order operator hedge in the regime where the
first-order hedge is itself accurate.
\item In the large-clock-shock regime ($\eta_{Y}=0.4$), HGB on
operator features outperforms the pure operator hedge by
$+0.14$ in $R^{2}$ ($0.65$ vs $0.51$): the predicted regime
where ML augments first-order analytic hedging with
higher-order corrections. The $\Gamma_{\mathrm{field}}$ and
$\Gamma_{\mathrm{clock}}$ squared features carry the
second-order content.
\item ML on global non-compositional features (set N) never
matches either ML+Op or the pure operator hedge in any regime
tested (except the very smooth $\varepsilon=2.4$ where it
reaches $0.84$ vs ML+Op's $0.97$). The structural feature
advantage $\Delta R^{2}\approx 0.5$--$0.7$ over Set N is
stable across payoffs.
\item Against the stronger state-conditioned baseline Set N+,
the gap shrinks to $0.35$--$0.49$ in $R^{2}$ at the main
regime---larger than a pure ``operator features as local state
encoding'' account would predict.
At very small $\eta_{Y}=0.02$ (deep Taylor regime) the gap
collapses to $0.06$, consistent with local-baseline approximation
of first-order shock structure; from $\eta_{Y}=0.05$ onward
ML+Op holds a $0.22$--$0.53$ margin over ML+N+. The structural
advantage is therefore not that the operator features are
state-conditioned, but that they are state-conditioned
\emph{in the directions} $f'(X)\,\delta\Bvar_{\varepsilon}(\Yclock)$
and
$f'(X)\,\partial_{x}\Bvar_{\varepsilon}(\Yclock)\,\delta\Yclock$
that the chain rule selects.
\item Individual operator features show heterogeneous
importance: for the quadratic payoff, two single features
carry $0.13$--$0.19$ $R^{2}$ each; for the other payoffs the
joint contribution dominates and individual leave-one-out
drops are small. ML benefits from the \emph{set} of operator
features rather than any one of them.
\end{itemize}

\paragraph{Verdict.}
Operator derivatives and ML are complementary, not competing.
The operator derivative supplies six state-conditioned features
that raise out-of-sample $R^{2}$ by $\sim 0.5$--$0.7$ over a
global non-compositional baseline and by $0.35$--$0.49$ over a
strictly stronger state-conditioned baseline (Set N+) that
matches the local PCA / nearby-grid / random-Fourier capacity
without the operator directions. In the mid-regime, a linear
Ridge regression on operator features reproduces the analytic
first-order operator hedge to within sampling error: the
operator hedge is optimally implementable by ML given the
right features. In the large-clock-shock regime where
first-order Taylor breaks down, gradient-boosted ML on
operator features (including the squared $\Gamma$ terms)
outperforms the pure analytic operator hedge by capturing
higher-order corrections that the first-order formula omits.
The framework's practical role is feature engineering for
compositional risk: identify compositional features
analytically---and specifically the chain-rule-selected
directions---then let ML do nonlinear combinations and
higher-order learning.

\medskip

The features and operator hedge are formulated on the mollified
field $\Bvar_{\varepsilon}$; the unmollified limit
$\varepsilon\downarrow 0$ is open and is not treated here. The
negative results ``ML+N $<$ ML+Op'' and ``ML+N+ $<$ ML+Op'' are
bounded by HGB on the fixed non-compositional and
state-conditioned baselines tested here. This does not rule out
every high-capacity raw-field architecture; it rules out the
tested non-compositional and state-conditioned baselines under a
controlled fixed-protocol comparison. A capacity-unbounded
architecture trained end-to-end on the raw scenery/clock inputs
might recover analogous representations internally, in which
case the framework's role would shift from feature construction
to feature validation, with the operator derivative identifying
which compositional features the learned representation must
recover to be optimal.

%==============================================================
\section*{Acknowledgements}
%==============================================================

AI tools were used as interactive assistants for drafting,
reorganizing, and generating mathematical arguments, under the
author's direction. The author is responsible for all
mathematical content and any errors.

%==============================================================
\begin{thebibliography}{99}

\bibitem{BidegreeAtlas}
R.~Fontes, \emph{A Bidegree Atlas for the Iterated-Brownian-Motion
  It\^o Correction}, manuscript (2026).

\bibitem{Burdzy}
K.~Burdzy, \emph{Some path properties of iterated Brownian motion},
  in: Seminar on Stochastic Processes, 1992, Progr.\ Probab.,
  vol.~33, Birkh\"auser, Boston, MA, 1993, pp.~67--87.

\bibitem{EnergyFirst}
R.~Fontes, \emph{Energy-First Stochastic Calculus: Iterated
  Brownian Motion as a Test Case}, manuscript (2026).

\bibitem{FCDC}
R.~Fontes, \emph{First-Chaos Dual Closure of the Iterated
  Brownian Motion It\^o Correction}, manuscript (2026).

\bibitem{FriedmanGBM}
J.~H. Friedman, \emph{Greedy function approximation: a gradient
  boosting machine}, Ann.\ Statist.\ \textbf{29} (2001),
  no.~5, 1189--1232.

\bibitem{GammaHierarchyDivergence}
R.~Fontes, \emph{From Carr\'e-du-Champ Hierarchies to Primitive
  Divergence}, manuscript (2026).

\bibitem{NualartMalliavin}
D.~Nualart, \emph{The Malliavin Calculus and Related Topics},
  2nd ed., Probability and Its Applications, Springer, Berlin,
  2006.

\bibitem{OperatorClarkOcone}
R.~Fontes, \emph{Operator Clark--Ocone Formula from Primitive
  Divergence}, manuscript (2026).

\end{thebibliography}

\end{document}
